% from ega4-9-fr.tex
\oldpage[IV-3]{66}

de (8.1.2, \emph{a)}), on peut (en remplaçant au besoin $S$ par un voisinage
de $\eta$) supposer que $h\in\Gamma(X_\eta,\mathcal O_{X_\eta})$ est de la
forme $g_\eta$, où $g\in\Gamma(X,\mathcal O_X)$~; soit
$u:\mathcal F\to\mathcal F$ l'endomorphisme de $\mathcal F$ défini par la
multiplication par $g$~; par hypothèse, on a $u_\eta=0$, donc (9.4.5)
l'endomorphisme $u_s:\mathcal F_s\to\mathcal F_s$ défini par la multiplication
par $g_s$, est nul dans un voisinage de $\eta$. D'autre part, soit
$v:\mathcal O_X\to\mathcal O_X$ l'endomorphisme défini par la multiplication
par $g$~; comme $B_\eta$ est intègre et $h=g_\eta\neq0$, $v_\eta$ est injectif,
et il résulte donc de (9.4.5) que $v_s$ est un endomorphisme injectif de
$\mathcal O_{X_s}$ pour $s$ voisin de $\eta$, autrement dit, $g_s$ est un
élément $\mathcal O_{X_s}$-régulier pour ces valeurs de $s$~; donc
$\mathcal F_s$ est de torsion dans un voisinage de $\eta$.

$2^\circ$ $\eta\in S-E$. Dire qu'un $B_\eta$-module de type fini $M_\eta$
n'est pas un module de torsion signifie que son quotient $M_\eta/T$ par son
sous-module de torsion est $\neq0$, et comme c'est un $B_\eta$-module sans
torsion de type fini, il est isomorphe à un sous-module d'un $B_\eta$-module
$B_\eta^n$~; il existe par suite un homomorphisme
$w:M_\eta\to B_\eta$ qui est $\neq0$. Appliquant (8.5.2, (i)) suivant la
méthode de (8.1.2, \emph{a)}), on en déduit (en remplaçant au besoin $S$ par un
voisinage de $\eta$) qu'il existe un homomorphisme
$v:\mathcal F\to\mathcal O_X$ tel que $v_\eta=\widetilde w$. L'hypothèse
$v_\eta\neq0$ entraîne donc (9.4.5) que $v_s\neq0$ dans un voisinage de
$\eta$, et comme $X_s$ est localement intègre, $\mathcal F_s$ n'est pas de
torsion pour ces valeurs de $s$.

$3^\circ$ $\eta\in E'$. Comme $M_\eta$ est un $B_\eta$-module de type fini
sans torsion, il existe un homomorphisme injectif
$w:M_\eta\to B_\eta^n$. Utilisant (8.5.2, (i)) et (9.4.5) comme dans le
$2^\circ$ (en restreignant au besoin $S$), on en déduit qu'il existe un
homomorphisme $v:\mathcal F\to\mathcal O_X^n$ tel que $v_\eta=\widetilde w$,
et que pour $s$ voisin de $\eta$,
$v_s:\mathcal F_s\to\mathcal O_{X_s}^n$ est injectif~; pour ces valeurs de
$s$, $\mathcal F_s$ est donc sans torsion.

$4^\circ$ $\eta\in S-E'$. Soit $T$ le sous-module de torsion de $M_\eta$~;
par hypothèse $T\neq0$, et $T$ est de type fini puisque $M_\eta$ est
noethérien. Utilisant cette fois (8.5.2, (i) et (ii)) on voit (en restreignant
au besoin $S$) qu'il existe un $\mathcal O_X$-Module cohérent $\mathcal G$ et
un homomorphisme injectif $u:\mathcal G\to\mathcal F$ tels que
$\mathcal G_\eta=\widetilde T$ et que $u_\eta$ soit l'injection canonique
$\widetilde T\to\mathcal F_\eta$. Il résulte alors du $1^\circ$ et de (1.8.6)
que dans un voisinage de $\eta$, $\mathcal G_s$ est un
$\mathcal O_{X_s}$-Module de torsion $\neq0$, et d'autre part il résulte de
(9.4.5) que dans un voisinage de $\eta$, $u_s$ est injectif. On en conclut que
dans un voisinage de $\eta$, le sous-Module de torsion de $\mathcal F_s$ n'est
pas nul. C.Q.F.D.

\begin{remark}[9.4.9]
\label{IV.9.4.9-fr}
La propriété «~$X$ est un $k$-préschéma algébrique localement intègre~» ne
vérifie pas la condition (9.2.1, (i)), et il n'est donc pas certain que
l'énoncé (9.4.8) soit encore valable quand on ne fait aucune hypothèse sur $S$
et qu'on suppose seulement que $f$ est un morphisme de présentation finie et
$\mathcal F$ un $\mathcal O_X$-Module de présentation finie. Considérons
toutefois le cas particulier suivant~: $S_0$ étant un préschéma localement
noethérien, soient $f_0:X_0\to S_0$ un morphisme de type fini, tel que les
fibres $(X_0)_{s_0}$ soient localement intègres (pour tout $s_0\in S_0$), et
$\mathcal F_0$ un $\mathcal O_{X_0}$-Module cohérent~; soit $g:S\to S_0$ un
morphisme quelconque, posons
$X=X_0\times_{S_0}S$, $\mathcal F=\mathcal F_0\otimes_{\mathcal O_{S_0}}
\mathcal O_S$, $f=f_0\times1_S:X\to S$, et \emph{supposons} que pour tout
$s\in S$, la fibre $X_s$ soit encore localement intègre. Alors l'ensemble $E$
(resp. $E'$) des $s\in S$ tels que $\mathcal F_s$ soit de torsion (resp. sans

